\section{Short literature pass on censored thermal observations}

\paragraph{Main takeaway.}
The censored-data likelihood itself is not the main novelty. Tobit models, censored regression, censored Gaussian processes, and inequality-constrained data assimilation are all established areas. The promising opening for this project is to bring these ideas into a thermal-camera/additive-manufacturing setting, where saturation is caused by a physical process and the saturated pixels can be used as inequality information.

Based on the papers reviewed, the two strongest directions are:
\begin{enumerate}
    \item data assimilation or inverse modeling with censored thermal observations;
    \item neural-network surrogate modeling trained directly with a censored or inequality-aware loss.
\end{enumerate}
Gaussian-process methods are still useful as a small-scale baseline or diagnostic, but they may not be the main contribution because censored/Tobit GP methods are already fairly developed.

\paragraph{Paper-by-paper takeaways.}
\begin{itemize}
    \item \textit{Amemiya, Tobit Models: A Survey.}
    This is useful for the classical Tobit framework and terminology. It makes clear that censored and truncated regression models are well-established, so the basic likelihood should be treated as background rather than novelty.

    \item \textit{Miller and Halpern, Regression with Censored Data.}
    This reviews older censored-regression methods, including survival-analysis-style approaches. It is useful background, but less directly connected to spatial thermal fields or physical inverse problems.

    \item \textit{Lai and Ying, Estimating a Distribution Function with Truncated and Censored Data.}
    This is a more theoretical statistics paper on estimation with truncated and censored data. It is useful for statistical depth, but not the most direct guide for the numerical experiments.

    \item \textit{Wani et al., Parameter Estimation of Hydrologic Models Using a Likelihood Function for Censored and Binary Observations.}
    This is a strong analogy for our problem. It shows that censored or even binary observations can still inform parameters of a physical model when they are included through a formal likelihood.

    \item \textit{Thacker, Data Assimilation with Inequality Constraints.}
    This supports the idea of incorporating inequality constraints directly into a variational data-assimilation framework, instead of assimilating first and correcting constraint violations afterward.

    \item \textit{Lauvernet et al., A Truncated Gaussian Filter for Data Assimilation with Inequality Constraints.}
    This extends the inequality-constraint idea to a truncated Gaussian filtering method. It shows that inequality information can be useful in a large physical-model setting, but the application is ocean modeling rather than thermal-camera saturation.

    \item \textit{Maatouk and Bay, Gaussian Process Emulators for Computer Experiments with Inequality Constraints.}
    This builds GP emulators satisfying constraints such as boundedness, monotonicity, and convexity. It is useful GP background, but the constraints are mostly on emulator behavior rather than on censored sensor observations.

    \item \textit{Basson et al., Variational Tobit Gaussian Process Regression.}
    This gives a sparse variational GP method for censored observations. It is close to a direct GP version of our observation model, which suggests that a pure censored-GP direction may not be novel enough unless combined with thermal physics or AM-specific structure.

    \item \textit{Chakraborty and Chakraborty, Efficient and Scalable Tobit Gaussian Process Regression.}
    This pushes Tobit GP regression toward scalable computation. It reinforces that scalable censored GP modeling is an active and already-developed area.

    \item \textit{Spiller et al., The Zero Problem.}
    This uses a zero-censored GP idea for range-constrained simulator outputs. The range-constraint viewpoint is relevant, but their setting concerns constrained simulator output rather than camera saturation as an observation process.

    \item \textit{Guo et al., Machine Learning for Metal Additive Manufacturing.}
    This motivates moving from black-box ML toward physics-informed data-driven modeling in metal AM. It supports using physically or statistically meaningful losses in surrogate models.

    \item \textit{Faegh et al., Review on Physics-Informed ML for Process-Structure-Property Modeling in AM.}
    This organizes PIML into physics-based features, architectures, and loss functions. Our censored or inequality-aware loss fits naturally into the loss-function category.
\end{itemize}

\paragraph{What seems mostly done.}
The classical Tobit/censored-likelihood framework is already established. Censored regression is mature, censored Gaussian-process regression has variational and scalable versions, and generic inequality-constrained GP emulators already exist. Physics-informed ML for AM is also a broad review-level topic. Therefore, the project should avoid framing the contribution as simply ``using a censored likelihood'' or ``doing a Tobit GP.'' These are better used as foundations or baselines.

\paragraph{Most promising direction 1: data assimilation or inverse modeling.}
This is the strongest physics-facing direction. The heat equation or a thermal simulator gives the physical model, while the camera provides censored observations:
\[
    Y_i = \min\{H_i(u) + \varepsilon_i, T_{\max}\},
    \qquad
    \varepsilon_i \sim \mathcal N(0,\sigma^2),
\]
where $u$ could denote the thermal state, heat-source parameters, material parameters, or boundary/input parameters. The observation loss could be written as
\[
    \mathcal L_{\mathrm{obs}}(u)
    =
    \sum_{i\in\mathcal U}
    \frac{(Y_i-H_i(u))^2}{2\sigma^2}
    -
    \sum_{i\in\mathcal S}
    \log \Phi\!\left(\frac{H_i(u)-T_{\max}}{\sigma}\right).
\]
This direction is promising because it asks whether saturated pixels improve inference of physically meaningful quantities, not just image reconstruction. A clean first study could infer heat-source parameters from clipped synthetic thermal data and compare exact, discard, hinge, and censored likelihood treatments.

\paragraph{Most promising direction 2: neural-network surrogate modeling with censored loss.}
This is the strongest ML-facing direction. Instead of reconstructing the missing high-temperature peak first, we can train a surrogate directly from censored thermal fields. For a network prediction $T_\eta(x;\theta)$, one possible inequality-aware loss is
\[
    \mathcal L_{\mathrm{hinge}}(\eta)
    =
    \sum_{i\in\mathcal U}
    \bigl(T_\eta(x_i;\theta)-Y_i\bigr)^2
    +
    \lambda
    \sum_{i\in\mathcal S}
    \bigl[T_{\max}-T_\eta(x_i;\theta)\bigr]_+^2.
\]
A probabilistic version would replace the hinge term with the censored Gaussian likelihood term. The clean experiment would compare three strategies: training on clipped fields as exact, reconstructing first and then training, and training directly with a censored or hinge loss. Since synthetic data gives the true uncensored fields, we can evaluate which strategy best predicts the full field and peak-temperature quantities on held-out examples.

\paragraph{Recommendation.}
Use GP methods as a small baseline or explanatory figure if useful, but center the next stage around data assimilation/inverse modeling and censored-loss surrogate modeling. Together, these give a balanced story: physical interpretability and uncertainty on one side, and practical ML surrogate training for additive manufacturing on the other.
